% LuaLaTeX document — compile with: lualatex cube_eshelby_explanation.tex
\documentclass[11pt,a4paper]{article}

\usepackage{fontspec}
\setmainfont{Latin Modern Roman}
\setmonofont{Latin Modern Mono}

\usepackage{amsmath,amssymb,bm}
\usepackage{booktabs}
\usepackage{geometry}
\geometry{margin=2.5cm}
\usepackage{hyperref}
\hypersetup{colorlinks=true,linkcolor=blue!60!black,citecolor=green!50!black,urlcolor=blue!70!black}
\usepackage{enumitem}
\usepackage{xcolor}

\newcommand{\D}{\Delta}
\newcommand{\order}{\mathcal{O}}
\newcommand{\bG}{\bm{G}}
\newcommand{\bT}{\bm{T}}
\newcommand{\bI}{\bm{I}}
\newcommand{\bC}{\bm{C}}
\newcommand{\bs}{\bm{\sigma}}
\newcommand{\be}{\bm{\varepsilon}}
\newcommand{\bu}{\bm{u}}

\title{Self-Consistent T-Matrix for Cubic Scatterers:\\
Static Eshelby, Dynamic Coupling, and the Unified Theory}
\author{Research Notes}
\date{March 2026}

\begin{document}
\maketitle

\tableofcontents
\vspace{1em}

%======================================================================
\section{The physical problem}
%======================================================================

We want to model seismic wave propagation through a medium filled with
cubic heterogeneities---a space-filling lattice of cubes, each with
different elastic properties than the background.  Every cube scatters
the wavefield, and the scattered field from each cube modifies the
field incident on every other cube.  The goal is an \emph{effective
medium}: replace the heterogeneous lattice with a homogeneous medium
that produces the same coherent wavefield.

This requires a \textbf{T-matrix} for each cube that is
\emph{self-consistent} in two senses simultaneously:
\begin{enumerate}[nosep]
  \item \textbf{Static self-consistency (Eshelby):} The internal field
    inside each cube is modified by the cube's own stiffness contrast.
    A stiff cube resists deformation; a soft cube amplifies it.  This
    is the classical Eshelby concentration factor.
  \item \textbf{Dynamic self-consistency (lattice coupling):} At finite
    frequency, each cube's scattered field propagates to all other cubes,
    modifying their incident fields.  The effective medium must account
    for this inter-site coupling.
\end{enumerate}

The central question of this document: \emph{how do we combine both
into a single self-consistent T-matrix for a cubic scatterer?}

%======================================================================
\section{What we have now: two separate pieces}
%======================================================================

\subsection{Piece 1: Static Eshelby amplification (the cube T-matrix)}

The existing module \texttt{effective\_contrasts.py} computes the
self-consistent T-matrix for a single cube embedded in a
\emph{background} medium.  The key integral is over the Green's tensor
second derivatives:
\begin{equation}\label{eq:Iijkl}
  I_{ijkl} = \int_{\text{cube}}
    \frac{\partial^2 G_{ij}(\bm{x})}{\partial x_k \partial x_l}\,dV,
\end{equation}
which decomposes by cubic symmetry into three coefficients:
\begin{equation}\label{eq:ABC}
  I_{ijkl} = A^c\,\delta_{ij}\delta_{kl}
            + B^c(\delta_{ik}\delta_{jl} + \delta_{il}\delta_{jk})
            + C^c\,E_{ijkl},
\end{equation}
where $E_{ijkl} = 1$ only when $i{=}j{=}k{=}l$.  The coefficient
$C^c$ is specific to the cube---it vanishes for a sphere.

These give rise to \textbf{four amplification factors}:
\begin{align}
  A_u &= \frac{1}{1 - \omega^2 \D\rho\,\Gamma_0}
    & &\text{(displacement/density),} \label{eq:amp_u}\\[4pt]
  A_\theta &= \frac{1}{1 - 3T_1^c - 2T_2^c - T_3^c}
    & &\text{(volumetric strain),} \label{eq:amp_theta}\\[4pt]
  A_e^{\text{off}} &= \frac{1}{1 - 2T_2^c}
    & &\text{(off-diagonal shear),} \label{eq:amp_eoff}\\[4pt]
  A_e^{\text{diag}} &= \frac{1}{1 - 2T_2^c - T_3^c}
    & &\text{(diagonal deviatoric),} \label{eq:amp_ediag}
\end{align}
where $T_1^c$, $T_2^c$, $T_3^c$ are coupling coefficients that depend
on $A^c$, $B^c$, $C^c$ and the material contrast $\D\lambda$, $\D\mu$.
These amplification factors already include
frequency dependence through the Taylor-expanded Green's tensor (the
integral \eqref{eq:Iijkl} is computed analytically at finite $\omega$).

For a sphere, $C^c = 0$ and the four reduce to three:
$A_e^{\text{off}} = A_e^{\text{diag}}$.  For a cube, the two shear
channels split, producing \textbf{cubic anisotropy} in the effective
shear modulus even from an isotropic inclusion.

\subsection{Piece 2: Dynamic lattice coupling (the lattice Green's tensor)}

The module \texttt{lattice\_greens.py} computes the inter-site
Green's tensor for a periodic lattice of scatterers.  The Foldy--Lax
system for $N$ sites reads
\begin{equation}\label{eq:foldylax}
  \bu_m^{\text{exc}} = \bu_m^{\text{inc}}
    + \sum_{n \neq m} \bG(\bm{x}_m - \bm{x}_n,\omega)\,
      \bT_n\,\bu_n^{\text{exc}},
\end{equation}
where $\bG$ is the elastodynamic Green's tensor and $\bT_n$ is the
T-matrix at site~$n$.

\subsection{The gap: these two pieces are independent}

Currently, the Eshelby amplification (Piece~1) is computed in the
\emph{background} medium.  The lattice coupling (Piece~2) uses the
\emph{background} Green's tensor.  Neither knows about the other:
\begin{itemize}[nosep]
  \item The Eshelby integral \eqref{eq:Iijkl} uses the background
    Green's tensor $\bG_0$, not the effective medium $\bG^*$.
  \item The lattice Foldy--Lax system \eqref{eq:foldylax} uses a
    T-matrix computed in the background, not in the effective medium.
\end{itemize}
At low contrast this is fine (the effective medium $\approx$ background).
At finite contrast, the two should be coupled: the effective medium
modifies both the Eshelby concentration factor \emph{and} the
inter-site propagation simultaneously.

%======================================================================
\section{What the literature says}
\label{sec:literature}
%======================================================================

\subsection{The sphere solution (Sabina \& Willis, 1988--1993)}

Sabina, Smyshlyaev \& Willis~\cite{sabina1993} solved the
simultaneous static\,+\,dynamic self-consistency problem for
\textbf{spherical} inclusions.  Their approach:
\begin{enumerate}[nosep]
  \item Embed a single sphere in an \emph{effective medium} with
    unknown moduli $\bC^*$ and wavenumber $k^*$.
  \item Solve the scattering problem for that sphere in the effective
    medium, obtaining a T-matrix $\bT(\bC^*, k^*, \omega)$.
  \item Demand self-consistency: the mean scattered field vanishes,
    $\langle\bT\rangle = 0$ (the coherent potential approximation, CPA).
  \item Iterate until $\bC^*$ converges.
\end{enumerate}
This \emph{automatically} produces the correct Eshelby factors in the
static limit (because the sphere in the effective medium has the
Eshelby concentration factor of the effective medium, not the
background) \emph{and} the correct dynamic dispersion at finite
frequency.  The two self-consistencies are not separate---they emerge
from the same Green's tensor $\bG^*$ of the effective medium.

\textbf{Limitation:} restricted to spherical (and later spheroidal)
inclusions, because the Eshelby tensor is uniform inside only for
ellipsoids.

\subsection{Dynamic Eshelby tensor for cubes (Wang et al., 2005; Kanaun \& Levin, 2008)}

Wang, Michelitsch, Gao \& Levin~\cite{wang2005} computed the
\textbf{dynamic Eshelby tensor} $\bm{S}(\bm{x}, \omega)$ for cubic
and prismatic inclusions using Helmholtz potentials.  Unlike the
static case (where the Eshelby tensor is non-uniform inside
non-ellipsoidal bodies), the dynamic Eshelby tensor is inherently
position-dependent even for ellipsoids.  What matters for the
T-matrix is the \emph{volume average}:
\begin{equation}\label{eq:Savg}
  \bar{\bm{S}}(\omega) = \frac{1}{V}\int_{\text{cube}}
    \bm{S}(\bm{x}, \omega)\,dV.
\end{equation}

Kanaun \& Levin~\cite{kanaun2008} then developed the
\textbf{effective field method} (EFM), a self-consistent scheme that
works with the volume-averaged quantities.  Their key insight: even
though $\bm{S}(\bm{x},\omega)$ varies inside the cube, the
self-consistent condition only needs the volume-averaged T-matrix, and
this can be computed for any shape.

\textbf{Verdict:} This is the closest existing framework to what we
need.  But nobody has applied it to a \emph{space-filling lattice}
of cubes with the full Eshelby delta-function handling.

\subsection{Willis (1981): the variational foundation}

Willis~\cite{willis1981} proved that the self-consistent scheme is
justified at all frequencies through a variational principle.  The
effective constitutive relation is \emph{nonlocal} (Willis materials:
stress depends on both strain and velocity), and the variational
principle applies for \textbf{arbitrary microstructure}---not just
spherical inclusions.  The self-consistent scheme is the saddle-point
approximation to the variational bounds.

\subsection{What nobody has done}

No paper in the open literature computes a \emph{simultaneously
static and dynamic self-consistent T-matrix for cubic (non-ellipsoidal)
scatterers on a lattice}.  The pieces exist:
\begin{itemize}[nosep]
  \item Sabina \& Willis show how to do it for spheres
  \item Kanaun \& Levin show how to compute the dynamic Eshelby
    tensor for cubes
  \item Willis's variational principle justifies the scheme for
    arbitrary shapes
\end{itemize}
But the synthesis---applying the Sabina--Willis self-consistent
iteration with the Kanaun--Levin cube Eshelby tensor on a periodic
lattice---has not been done.

%======================================================================
\section{The unified theory: what needs to change}
\label{sec:unified}
%======================================================================

\subsection{The key insight: one Green's tensor does everything}

The Eshelby concentration factor and the dynamic inter-site coupling
are \emph{not} two separate physics.  They are two regimes of the
\emph{same} Green's tensor integral.

The Green's tensor of the effective medium $G^*_{ij}(\bm{x}, \omega)$
has three parts:
\begin{equation}\label{eq:Gdecomp}
  G^*_{ij}(\bm{x}) = \underbrace{G^{\text{Kelvin}}_{ij}(\bm{x})}_{1/r\text{ static}}
    + \underbrace{G^{\text{smooth}}_{ij}(\bm{x}, \omega)}_{\text{dynamic radiation}}
\end{equation}
When we integrate $\partial^2 G^*/\partial x_k \partial x_l$ over the
cube volume:
\begin{itemize}[nosep]
  \item The \textbf{Kelvin part} produces the $1/r^3$ Eshelby
    singularity at $r=0$ --- this gives the static
    depolarization tensor ($A^c$, $B^c$, $C^c$).
  \item The \textbf{smooth part} produces $\order(ka)^2$ frequency
    corrections --- this is the dynamic radiation coupling.
\end{itemize}
Both live in the integral \eqref{eq:Iijkl}.  Currently we compute
this integral using the \emph{background} Green's tensor.  The
correction is to use the \emph{effective medium} Green's tensor
instead.

\subsection{The self-consistent iteration}

Following Sabina--Willis~\cite{sabina1993}, adapted to cubes:

\begin{enumerate}
  \item \textbf{Initialize} $\bC^* = \bC_0$ (background moduli),
    $k^* = k_0$ (background wavenumber).

  \item \textbf{Compute the cube T-matrix in the effective medium.}
    Use $\bG^*(\bm{x}, \omega)$ (with moduli $\bC^*$ and wavenumber
    $k^*$) in all integrals:
    \begin{itemize}[nosep]
      \item $\Gamma_0^* = \int_{\text{cube}} G^*_{ii}(\bm{x})\,dV$
        \quad (displacement self-energy)
      \item $A^{c*}, B^{c*}, C^{c*}$ from $\int_{\text{cube}}
        \partial^2 G^*_{ij}/\partial x_k \partial x_l\,dV$
      \item The Eshelby delta-function part comes from the static
        Kelvin component of $\bG^*$ (with $\mu^*$, $\nu^*$).
      \item The smooth part comes from the radiation component
        of $\bG^*$ (with $k^*_P$, $k^*_S$).
    \end{itemize}
    This produces amplification factors
    $A_u^*$, $A_\theta^*$, $A_e^{*\text{off}}$, $A_e^{*\text{diag}}$
    and effective contrasts
    $\D\rho^*$, $\D\lambda^*$, $\D\mu^*_{\text{off}}$,
    $\D\mu^*_{\text{diag}}$ --- all relative to $\bC^*$, not $\bC_0$.

  \item \textbf{Update the effective medium.}  The CPA condition
    $\langle\bT\rangle = 0$ requires the effective medium to be the
    one in which the average scattering vanishes.  For a two-phase
    composite (volume fraction $\phi$ of inclusions in matrix):
    \begin{equation}\label{eq:CPA}
      \bC^*_{\text{new}} = \bC_0 + \phi\,\D\bC\,\bm{A}^*,
    \end{equation}
    where $\bm{A}^*$ is the concentration tensor built from
    $A_\theta^*$, $A_e^{*\text{off}}$, $A_e^{*\text{diag}}$.
    For a space-filling lattice ($\phi = 1$), this simplifies.

    The effective wavenumber updates via the Foldy--Lax dispersion
    relation:
    \begin{equation}\label{eq:kstar}
      (k^*)^2 = k_0^2 + n_0\,\bT_{\text{forward}}(\bC^*, k^*),
    \end{equation}
    where $n_0$ is the number density and $\bT_{\text{forward}}$ is
    the forward-scattering amplitude.

  \item \textbf{Iterate} Steps 2--3 until $\bC^*$ and $k^*$ converge.
\end{enumerate}

\subsection{What changes in the code}

The existing function \texttt{compute\_cube\_tmatrix(omega, a, ref,
contrast)} takes a \texttt{ReferenceMedium} (the background).  In the
unified scheme:
\begin{itemize}[nosep]
  \item The \texttt{ref} argument becomes the \emph{effective medium}
    $(\alpha^*, \beta^*, \rho^*)$, which changes at each iteration.
  \item The \texttt{contrast} becomes $\D\bC = \bC_{\text{inclusion}}
    - \bC^*$ (the contrast relative to the effective medium, not the
    background).
  \item The Eshelby integrals ($A^c$, $B^c$, $C^c$) and Green's
    tensor integral ($\Gamma_0$) are computed using the effective
    medium's velocities $\alpha^*$, $\beta^*$, not $\alpha_0$,
    $\beta_0$.
\end{itemize}
The key point: \textbf{the existing code already does the right
calculation}.  We just need to call it with the effective medium
instead of the background, and wrap it in an iteration loop.

%======================================================================
\section{What the current module does (and its limitations)}
\label{sec:current}
%======================================================================

The new module \texttt{cube\_eshelby.py} computes the four
concentration factors by comparing the full self-consistent T-matrix
against the Born (linearized) limit:
\begin{equation}
  E_{\text{channel}} = \frac{\D X^*_{\text{full}}}{\D X^*_{\text{Born}}},
\end{equation}
where $\D X^*$ is the effective contrast for each channel.

\subsection{What it gets right}

\begin{itemize}[nosep]
  \item The four static amplification factors
    (Eqs.~\ref{eq:amp_u}--\ref{eq:amp_ediag}) are exact in the
    Rayleigh regime ($ka < 0.3$).
  \item The Born limit is computed cleanly by epsilon-scaling the
    full solver.
  \item The concentration ratios $E$ correctly quantify the
    departure from Born.
  \item The cubic anisotropy
    ($\D\mu^*_{\text{diag}} \neq \D\mu^*_{\text{off}}$) is captured.
\end{itemize}

\subsection{What it misses}

\begin{itemize}[nosep]
  \item \textbf{No effective-medium iteration.}  The Eshelby integrals
    use the background medium, not the effective medium.  This means
    the static self-consistency is relative to the background, not
    the true effective medium of the lattice.
  \item \textbf{No lattice dispersion.}  The amplification factors
    don't know about the lattice structure or the effective wavenumber.
  \item \textbf{Resonance extraction is approximate.}  At $ka > 0.3$,
    the 9$\times$9 composite T-matrix from the Foldy--Lax sub-cell
    solver does not decompose cleanly into 4~scalar concentration
    factors (there are off-diagonal couplings).
\end{itemize}

\subsection{The path forward}

The unified scheme (Section~\ref{sec:unified}) replaces the
background with the effective medium in \texttt{compute\_cube\_tmatrix},
adds the CPA iteration (Eq.~\ref{eq:CPA}), and couples to the
lattice dispersion (Eq.~\ref{eq:kstar}).  The existing code is
the inner kernel of this iteration---no new physics is needed,
only the outer loop.

%======================================================================
\section{Numerical results (current module)}
\label{sec:results}
%======================================================================

Using the standard test parameters ($\alpha_0 = 5000$~m/s,
$\beta_0 = 3000$~m/s, $\rho_0 = 2500$~kg/m$^3$) at $ka = 0.01$
(deep Rayleigh regime), with proportional contrast
$\D\lambda = \lambda_0 f$, $\D\mu = \mu_0 f$, $\D\rho = \rho_0 f$:

\begin{center}
\begin{tabular}{rcccc}
\toprule
\textbf{Contrast} & $\bm{A_\theta}$ & $\bm{A_e^{\textbf{off}}}$
  & $\bm{A_e^{\textbf{diag}}}$
  & $\bm{\D\mu^*_{\textbf{diag}} - \D\mu^*_{\textbf{off}}}$ \\
\midrule
  1\% & 0.9948 & 0.9957 & 0.9941 & $-3.6 \times 10^5$~Pa\\
  5\% & 0.9747 & 0.9789 & 0.9712 & $-8.6 \times 10^6$~Pa\\
 10\% & 0.9506 & 0.9586 & 0.9440 & $-3.3 \times 10^7$~Pa\\
 20\% & 0.9058 & 0.9206 & 0.8940 & $-1.2 \times 10^8$~Pa\\
 40\% & 0.8278 & 0.8528 & 0.8083 & $-4.0 \times 10^8$~Pa\\
\bottomrule
\end{tabular}
\end{center}

All factors are $< 1$ for stiff inclusions (the cube resists
deformation).  The density channel $A_u \approx 1.000$ throughout
(density has no static Eshelby correction; the deviation is
$\order(\omega^2)$).

\subsection{Key observations}

\begin{enumerate}[label=(\roman*)]

\item \textbf{Cube $A_\theta$ matches sphere $E_0$ exactly.}
At the test contrast ($\D\lambda = 2$~GPa, $\D\mu = 1$~GPa):
cube $A_\theta = 0.9591$, sphere
$K_0/(K_0 + \alpha_E \D K) = 0.9591$.  This is because the
volumetric part of the Eshelby tensor ($J_a$) is
\emph{universal} for any convex body:
$J_{a,kl} = -\tfrac{4\pi}{3}\,\delta_{kl}$
(from $\nabla^2(1/r) = -4\pi\delta$).

\item \textbf{Cube $A_e^{\text{off}}$ differs from sphere $E_2$.}
Sphere: $E_2 = 0.978$.  Cube: $A_e^{\text{off}} = 0.981$.
The difference arises from the geometry-specific surface integrals
($j_1$, $j_2$, $k_1$).

\item \textbf{Cubic anisotropy is real.}
$A_e^{\text{diag}} < A_e^{\text{off}}$: the cube scatters
diagonal deviatoric strain more strongly than off-diagonal shear.
The anisotropy grows quadratically with contrast.

\item \textbf{Scale-free.}
$A_\theta(a{=}0.5) = A_\theta(a{=}2.0) = 0.959079$ to six digits.
The factors depend on $ka$ and contrast ratios only.

\item \textbf{Born recovered at weak contrast.}  At 0.01\% contrast,
all factors are within $6 \times 10^{-5}$ of unity.

\end{enumerate}

%======================================================================
\section{The big picture: where this fits}
%======================================================================

The full self-consistent effective medium for a lattice of cubes
requires three nested calculations:

\begin{center}
\renewcommand{\arraystretch}{1.3}
\begin{tabular}{llp{7.5cm}}
\toprule
\textbf{Level} & \textbf{What} & \textbf{Status} \\
\midrule
Inner &
  Single-cube T-matrix in a given medium &
  Done: \texttt{compute\_cube\_tmatrix}.
  Takes any reference medium; computes $A^c$, $B^c$, $C^c$,
  $\Gamma_0$, and all four amplification factors analytically. \\
Middle &
  CPA iteration for effective moduli $\bC^*$ &
  Done: \texttt{cpa\_iteration.py}.  Calls the inner kernel
  repeatedly with updated $\bC^*$ until $\langle\bT\rangle = 0$.
  Converges in $\sim$5 iterations at moderate contrast. \\
Outer &
  Lattice dispersion $k^*(\omega)$ &
  Partially done: \texttt{lattice\_greens.py} provides the
  lattice Green's tensor.  Needs coupling to the CPA iteration
  to get the self-consistent effective wavenumber. \\
\bottomrule
\end{tabular}
\end{center}

The \texttt{cube\_eshelby.py} module characterises the
\emph{inner} kernel.  The \texttt{cpa\_iteration.py} module
implements the middle loop, described next.

%======================================================================
\section{The CPA iteration: implementation and subtlety}
\label{sec:cpa}
%======================================================================

\subsection{Algorithm}

The module \texttt{cpa\_iteration.py} implements the CPA
self-consistent loop:

\begin{enumerate}
  \item \textbf{Input:} A list of material phases
    $(\lambda_n, \mu_n, \rho_n, \phi_n)$ with $\sum_n \phi_n = 1$.
  \item \textbf{Initialize:} $\bC^* = \bC_{\text{Voigt}}
    = \sum_n \phi_n \bC_n$ (arithmetic average).
  \item \textbf{Iterate:}
    \begin{enumerate}[label=(\alph*)]
      \item Convert $\bC^*$ to an isotropic \texttt{ReferenceMedium}
        $(\alpha^*, \beta^*, \rho^*)$ using $\mu^* = \mu^*_{\text{off}}$.
      \item For each phase $n$, compute the contrast
        $\D\bC_n = \bC_n - \bC^*$ and run
        \texttt{compute\_cube\_tmatrix} to get the T-matrix effective
        contrasts $\D\rho^*_n$, $\D\lambda^*_n$, $\D\mu^{*\text{off}}_n$.
      \item Update: $\bC^* \leftarrow \bC^* + \sum_n \phi_n \D\bC^*_n$.
      \item Check convergence on the three isotropic channels.
    \end{enumerate}
  \item \textbf{Output:} Converged $\bC^*$ with cubic symmetry:
    $\lambda^*$, $\mu^*_{\text{off}}$, $\mu^*_{\text{diag}}$, $\rho^*$.
\end{enumerate}

\subsection{The cubic anisotropy subtlety}

The cube T-matrix produces \textbf{four} independent effective
contrasts ($\D\lambda^*$, $\D\mu^*_{\text{off}}$,
$\D\mu^*_{\text{diag}}$, $\D\rho^*$), but \texttt{compute\_cube\_tmatrix}
takes an \emph{isotropic} reference medium with only
\textbf{three} elastic parameters ($\lambda^*$, $\mu^*$, $\rho^*$).

This means we can satisfy only three CPA equations:
\begin{equation}
  \langle \D\rho^* \rangle = 0, \quad
  \langle \D\lambda^* \rangle = 0, \quad
  \langle \D\mu^{*\text{off}} \rangle = 0.
\end{equation}
The fourth equation $\langle \D\mu^{*\text{diag}} \rangle = 0$
cannot be independently enforced because the isotropic reference
cannot represent the cubic anisotropy of $\bC^*$.

The solution: iterate on the three isotropic channels, and derive
$\mu^*_{\text{diag}}$ as a \emph{diagnostic} after convergence.
The cubic anisotropy $\mu^*_{\text{diag}} - \mu^*_{\text{off}}$
arises from the $C^c$ coefficient (the geometric cubic term in the
Eshelby tensor) and is accumulated during the iteration.

At moderate contrast ($\sim$10\%), the cubic anisotropy
$|\mu^*_{\text{diag}} - \mu^*_{\text{off}}|$ is typically
$\order(10^{-4})$ relative to $\mu^*$---small but nonzero.
To achieve a fully self-consistent cubic effective medium
would require generalising \texttt{compute\_cube\_tmatrix} to accept
a cubic (not just isotropic) reference.

\subsection{Convergence properties}

\begin{itemize}[nosep]
  \item \textbf{Weak contrast:} converges in 2--3 iterations to
    tolerance $10^{-10}$; the result is close to the Voigt average.
  \item \textbf{Moderate contrast ($\sim$10\%):} converges in 5--10
    iterations.  Under-relaxation ($0 < \text{damping} < 1$)
    available for stability at high contrast.
  \item \textbf{Single-phase:} converges in 1 iteration to the
    input (a sanity check: if the medium is homogeneous, $\langle\bT\rangle = 0$
    trivially).
  \item \textbf{Effective medium is physical:} all moduli remain
    positive and bounded between the constituent phases.
\end{itemize}

%======================================================================
\section{Kennett layer stacking: the outer layer}
\label{sec:kennett}
%======================================================================

The module \texttt{kennett\_layers.py} implements the ``outer layer''
in the three-level architecture: it propagates coherent plane waves
vertically through a stack of CPA-homogenized layers using Kennett's
recursive reflection/transmission algorithm~\cite{kennett1983}.

\subsection{Algorithm}

The medium is divided into $N$ horizontal layers, each one cube thick
($d = 2a$).  Each layer is independently CPA-homogenized to obtain
isotropic effective properties $(\alpha^*_n, \beta^*_n, \rho^*_n)$.
The Kennett recursion then computes the cumulative downgoing
reflection matrix $\bm{R}_D$ by sweeping upward from the half-space:

\begin{enumerate}
  \item \textbf{Initialize:} $\bm{R}^{(N)}_D = 0$ (radiation condition
    at the half-space).
  \item \textbf{At each interface} $n = N{-}1, \ldots, 1$:
    \begin{enumerate}[label=(\alph*)]
      \item Compute the phase-shifted cumulative reflection:
        $\bm{M}_T = \bm{E}_n \bm{R}^{(n+1)}_D \bm{E}_n$,
        where $\bm{E}_n = \text{diag}(e^{i\omega\eta_n h_n},\,
        e^{i\omega\zeta_n h_n})$ is the one-way phase through layer~$n$.
      \item Compute the reverberation operator:
        $\bm{U} = (\bm{I} - \bm{R}_u \bm{M}_T)^{-1}$.
      \item Update:
        $\bm{R}^{(n)}_D = \bm{R}_d + \bm{T}_u \bm{M}_T \bm{U} \bm{T}_d$.
    \end{enumerate}
  \item \textbf{Output:} $\bm{R}^{(1)}_D(\omega, p)$, the $2 \times 2$
    P-SV reflection matrix at the top of the stack.
\end{enumerate}

Here $\bm{R}_d$, $\bm{R}_u$, $\bm{T}_d$, $\bm{T}_u$ are the
modified interfacial scattering coefficients (Aki \& Richards, 1980)
with $\sqrt{\eta\rho}$ normalization that ensures $\bm{T}_u = \bm{T}_d^T$
(reciprocity) and $\bm{R}_d^\dagger\bm{R}_d + \bm{T}_d^\dagger\bm{T}_d
= \bm{I}$ (energy conservation).

The SH system decouples into a scalar recursion with impedance
$Z_{\text{SH}} = \rho\beta^2\zeta$.

\subsection{CPA bridge}

The function \texttt{cubic\_to\_isotropic\_layer} converts a
\texttt{CubicEffectiveMedium} to an \texttt{IsotropicLayer} using
$\mu^*_{\text{off}}$ for the isotropic shear modulus (consistent with
\texttt{as\_reference\_medium}).  The cubic anisotropy
$\mu^*_{\text{diag}} - \mu^*_{\text{off}}$ is $\order(10^{-4})$
relative and is neglected in the isotropic Kennett propagation.

The function \texttt{cpa\_stack\_from\_phases} automates the full
pipeline: given per-layer phase lists, it runs \texttt{compute\_cpa}
on each layer and assembles the resulting \texttt{LayerStack}.

\subsection{Connection to the three-level architecture}

\begin{center}
\renewcommand{\arraystretch}{1.3}
\begin{tabular}{llp{7.5cm}}
\toprule
\textbf{Level} & \textbf{Module} & \textbf{Role} \\
\midrule
Inner &
  \texttt{effective\_contrasts.py} &
  Single-cube T-matrix in a given reference medium. \\
Middle &
  \texttt{cpa\_iteration.py} &
  Self-consistent effective moduli $\bC^*$ per layer
  (CPA: $\langle\bT\rangle = 0$). \\
Outer &
  \texttt{kennett\_layers.py} &
  Coherent plane-wave propagation through the
  CPA-homogenized layer stack (Kennett recursion). \\
\bottomrule
\end{tabular}
\end{center}

%======================================================================
\section{References}
%======================================================================

\begin{thebibliography}{99}

\bibitem{sabina1993}
F.J.\ Sabina, V.P.\ Smyshlyaev, J.R.\ Willis,
``Self-consistent analysis of waves in a matrix-inclusion composite,''
\emph{J.\ Mech.\ Phys.\ Solids} \textbf{41}, 1573--1588 (1993).

\bibitem{sabina1988}
F.J.\ Sabina, J.R.\ Willis,
``A simple self-consistent analysis of wave propagation in particulate
composites,''
\emph{Wave Motion} \textbf{10}, 127--142 (1988).

\bibitem{wang2005}
J.\ Wang, T.M.\ Michelitsch, H.\ Gao, V.M.\ Levin,
``On the solution of the dynamic Eshelby problem for inclusions of
various shapes,''
\emph{Int.\ J.\ Solids Struct.}\ \textbf{42}, 353--363 (2005).

\bibitem{kanaun2008}
S.\ Kanaun, V.\ Levin,
\emph{Self-Consistent Methods for Composites: Vol.~2.\ Wave
Propagation in Heterogeneous Materials},
Springer (2008).

\bibitem{willis1981}
J.R.\ Willis,
``Variational principles for dynamic problems for inhomogeneous
elastic media,''
\emph{Wave Motion} \textbf{3}, 1--11 (1981).

\bibitem{willis1977}
J.R.\ Willis,
``Bounds and self-consistent estimates for the overall moduli of
anisotropic composites,''
\emph{J.\ Mech.\ Phys.\ Solids} \textbf{25}, 185--202 (1977).

\bibitem{gubernatis1977}
J.E.\ Gubernatis, E.\ Domany, J.A.\ Krumhansl, M.\ Huberman,
``The Born approximation in the theory of the scattering of elastic
waves by flaws,''
\emph{J.\ Appl.\ Phys.}\ \textbf{48}, 2804--2811 (1977).

\bibitem{stanke1984}
F.E.\ Stanke, G.S.\ Kino,
``A unified theory for elastic wave propagation in polycrystalline
materials,''
\emph{J.\ Acoust.\ Soc.\ Am.}\ \textbf{75}, 665--681 (1984).

\bibitem{varadan1985}
V.K.\ Varadan, V.V.\ Varadan,
``Multiple scattering theory for wave propagation in discrete random
media,''
\emph{Int.\ J.\ Eng.\ Sci.}\ \textbf{22}, 1165--1176 (1985).

\bibitem{kim2020}
S.\ Kim, S.\ Torquato,
``Effective elastic wave characteristics of composite media,''
\emph{New J.\ Phys.}\ \textbf{22}, 123050 (2020).
arXiv:2012.15214.

\bibitem{parnell2008}
W.J.\ Parnell, I.D.\ Abrahams,
``A new integral equation approach to elastodynamic homogenization,''
\emph{Proc.\ R.\ Soc.\ A} \textbf{464}, 1461--1482 (2008).

\bibitem{buryachenko2001}
V.A.\ Buryachenko,
``Multiparticle effective field and related methods in micromechanics
of random structure composites,''
\emph{Appl.\ Mech.\ Rev.}\ \textbf{54}, 1--47 (2001).

\bibitem{kennett1983}
B.L.N.\ Kennett,
\emph{Seismic Wave Propagation in Stratified Media},
Cambridge University Press (1983).

\end{thebibliography}

\end{document}
